\chapter{Solução proposta}

Durante a primeira etapa do projeto, foi desenvolvida uma aplicação para validar a metodologia do trabalho. Essa aplicação utilizava o Java Microbenchmarking Harness para realizar o \textit{warmup} dos \textit{benchmarks} com o objetivo de coletar as métricas somente após as otimizações realizadas pela Java Virtual Machine (JVM). Além disso, foi utilizado o JRAPL \cite{jrapl} para medir o consumo de energia de cada iteração dos \textit{benchmarks}. Utilizando esse conjunto de ferramentas, realizou-se a medição de energia para três algoritmos simples de ordenação: Merge Sort, Quick Sort e Bubble Sort. Dessa forma, foi possível confirmar a viabilidade da metodologia, possibilitando a elaboração de testes mais complexos com outros tipos de algoritmos.

No início do PGC II, elaborou-se o plano de integrar à aplicação desenvolvida previamente os algoritmos do The Computer Language Benchmarks Game (CLBG) \cite{benchmarksgame}. Como demonstrado por \citeonline{programming_lang_energy}, o CLBG apresenta um conjunto robusto de algoritmos que podem ser usados para comparar, dentre outras métricas, o consumo de energia. Para analisar os algoritmos do CLBG foram selecionados a princípio três algoritmos diferentes: Fannkuchredux, Mandelbrot e Nbody. Apesar do resultado satisfatório, notou-se que os algoritmos do CLBG não possuem a diversidade e complexidade necessária para testar cargas de trabalho variadas. Além disso, as métricas disponibilizadas pelo \textit{benchmark} são limitadas, focando principalmente no tempo de execução. Isso tornava mais difícil um análise aprofundada do consumo de energia da JVM, já que havia poucas variáveis para comparar os \textit{benchmarks}. Isso motivou a mudança da metodologia para uma ferramenta mais robusta de \textit{benchmarking}, que é caso do DaCapo Chopin \cite{dacapochopin}.

% \begin{figure}
%    \centering
%    \includegraphics[scale=.4]{figs/energy_fannkuchredux_v1.png}
%    \caption{Consumo de energia do \textit{benchmark} Fannkuchredux.}
%    \label{fig:fannkuchredux}
% \end{figure}

% \begin{figure}
%    \centering
%    \includegraphics[scale=.4]{figs/energy_mandelbrot_v1.png}
%    \caption{Consumo de energia do \textit{benchmark} Mandelbrot.}
%    \label{fig:mandelbrot}
% \end{figure}

% \begin{figure}
%    \centering
%    \includegraphics[scale=.4]{figs/energy_nbody_v1.png}
%    \caption{Consumo de energia do \textit{benchmark} Nbody.}
%    \label{fig:nbody}
% \end{figure}

O DaCapo Chopin \cite{dacapochopin} contém 22 \textit{benchmarks}. Utilizando o método de \textit{principal component analysis} (PCA), os autores demonstraram que as cargas de trabalho disponibilizadas são variadas, ou seja, possuem comportamentos diferentes entre si. Isso possibilita observar o comportamento do consumo de energia para aplicações com perfis diversos. Sendo assim, foi desenvolvida uma extensão para o pacote DaCapo que possibilita medir o consumo de energia. Para isso, mais uma vez, optou-se pelo JRAPL \cite{jrapl} pelas razões já apresentadas. Utilizando essa extensão, foi possível testar 16 dos 22 \textit{benchmarks}, devido a problemas de dependências, os 6 \textit{benchmarks} restantes serão analisados na terceira parte do trabalho. Os resultados desses testes são discutidos na seção  \ref{chap:resultados}.